\section{Conclusion}

This project presented the underlying development of the prototype of
a modular digital audio workstation, implemented in C++, and is
currently under process as of publishing this report. The experimental
results using theoretical estimation indicate that all these implemented
audio components should potentially satisfy the real-time processing 
requirements. The mockup DAW session achieved a 
mean execution time of \SI{1039.295}{\micro\second} and a worst-case 
execution time of \SI{2245.130}{\micro\second}, remaining well below 
the \SI{11.61}{\milli\second} processing deadline for a 512 floating-point
sample audio buffer at a sampling rate of 44.1 kHz. These findings indicate 
that the current implementation is potentially capable of real-time audio processing.
Moreover, the total execution time of \SI{1186.609}{\micro\second} using a single thread
also indicates that it does not require rigourous multi-threading with just the small amount
of tracks, i.e., 6 tracks.

Beyond benchmark results, these implementations are also evaluated through
perceptual listening tests and signal analysis. All the Bode Plots and time-domain 
plots demonstrated the expected behaviour of the audio effects implemented, while
testing them individually. These effects are proven to show consistency with the 
expected results, which includes filtering, dynamic compression, reverb, etc. However, unison
synthesizer exhibited some audible beating effect caused by different frequencies 
between voices of the same amplitude, requiring some level of improvement.

These conclusions are currently limited to the estimations with specified components
and a mockup of a DAW session with a dummy list of instruments and audio effects only. And 
since the project is still under development, 
the measured performance does not represent the actual stress test of a complete DAW 
operating with a full live production session containing additional tracks, routing, 
automation, and user interaction. Further evaluation will be required as more 
functionalities are integrated into the system.

\section{Recommendations and Future Work}

It is highly recommended that a professional music project needs to be produced before testing them carefully, 
either by measuring callback execution time, or by monitoring CPU and RAM usage offered by the operating system, which is a 
better approach.
It is also recommended that the software is tested under various computer hardware specification, from lower-grade to
higher-grade.

Since software is not a physical resource, it can be created, deleted or modified without any additional cost 
(e.g. cost of materials).
Therefore they always keep evolving. When new or better solutions are found through enormous amount of research, 
they are either added and/or can replace currently existing solution, which could promise better results.

The digital audio workstation (DAW) is a very large-scale project, and it takes months or years of research and
development to complete the project as a solo developer. These kinds of large-scale software are often developed
by a team of programmers, audio designers, and DSP specialists, with their own area of expertise so that the job
can be done much faster using division of labour.
Which is why for starters, joining a community of audio developers (e.g. UVic CMCU),
attending conferences, and working on a project with them is highly encouraged. It even bring in to mentorship and guidance which
can go beyond primary/peer-reviewed sources, it order to achieve better and more efficient outcome.

In the meantime, some additional features can be added to this project. For e.g., 
ADSR (Attack, Decay, Sustain, Release), LFO (Low-frequency oscillation), compressor
sidechaining, stereo support, Arpeggiator, etc. Also that the Qt GUI is yet to be implemented, since every
DAW is equipped with a user interface that allows producers
to intuitively control parameters, display waveforms, drum rack for beat patterns, FFT-based spectrum analyzer, etc.,
Also, the sample generators currently generate one sample
at a time, and PortAudio accepts blocks of samples. Therefore, copying a block of memory
instead of storing them to a buffer sample-by-sample might improve the performance to a much
greater extent.

And finally, if a multi-threadeded architecure is employed to the DAW, it would be a good practice to implement 
Consumer-Producer Queue \cite{susser_adc2025} to prevent such resource contention.


% - Adding improvements previously mentioned in the Discussion section, 
% such as sending blocks directly to the audio engine instead of going sample-by-sample
% Supporting stereo sounds to mitigate beating effects
% ADSR, LFO, compressor sidechains, 

